\documentclass[a4paper,11pt]{article}

\usepackage{fontspec}
\setmainfont{Carlito}
\usepackage[margin=0.6in, top=0.5in, bottom=0.5in]{geometry}
\usepackage{titlesec}
\usepackage{enumitem}
\usepackage{hyperref}
\usepackage{xcolor}
\usepackage{microtype}
\usepackage{needspace}

\hyphenpenalty=10000 \exhyphenpenalty=10000 \emergencystretch=3em

\definecolor{headernavy}{RGB}{12, 83, 166}
\definecolor{linkblue}{RGB}{0, 102, 204}
\definecolor{muted}{RGB}{96, 102, 110}

\hypersetup{colorlinks=true, linkcolor=linkblue, urlcolor=linkblue}

\titleformat{\section}
  {\color{headernavy}\normalsize\bfseries}
  {}{0em}{\MakeUppercase}[\color{headernavy}\titlerule]

\titlespacing{\section}{0pt}{8pt}{4pt}

\setlist[itemize]{noitemsep, topsep=1pt, leftmargin=1.3em, partopsep=0pt, parsep=1pt}

\pagestyle{empty}
\setlength{\parindent}{0pt}
\linespread{1.02}

% ---- building blocks: bold = organisation/project, italic = role/stack, grey = dates/places ----
\newcommand{\role}[4]{% org, location, dates, title
  \textbf{#1}{\color{muted}\,\textperiodcentered\,\textit{#2}}\hfill{\color{muted}#3}\\
  \textit{#4}\par}
\newcommand{\roleNoLoc}[3]{\textbf{#1}\hfill{\color{muted}#2}\\ \textit{#3}\par}
\newcommand{\project}[3]{% name, stack, right-hand tag
  \textbf{#1}\hfill#3\\ {\color{muted}\small\textit{#2}}\par}
\newcommand{\gh}[1]{\href{#1}{\small GitHub}}
\newcommand{\badge}[1]{{\color{headernavy}\small\textbf{#1}}}
\newcommand{\skillgroup}[2]{{\color{headernavy}\textbf{#1}}\\ #2\par\vspace{3pt}}
\newcommand{\gap}{\vspace{5pt}}
\newcommand{\num}[1]{\textbf{#1}}   % result numbers only


\begin{document}

% ================= HEADER =================
\begin{center}
  {\LARGE\bfseries Thiruvel Andagurunathan Pandian} \\[3pt]
  {\color{headernavy}\small\textbf{Data Scientist}} \\[3pt]
  {\small
  Bristol, UK\nobreak\ \,\textbar\, \href{mailto:apthiruvel@gmail.com}{apthiruvel@gmail.com}\nobreak\ \,\textbar\, \href{tel:+447741016954}{+44 7741 016954}\nobreak\ \,\textbar\, \href{https://www.linkedin.com/in/thiruvel-a-p/}{https://www.linkedin.com/in/thiruvel-a-p/}\nobreak\ \,\textbar\, \href{https://github.com/Thiruvel-AP}{https://github.com/Thiruvel-AP}\nobreak\ \,\textbar\, \href{https://thiruvel-ap.github.io/PersonalPortfolio/}{https://thiruvel-ap.github.io/PersonalPortfolio/}}\par\vspace{4pt}
  \noindent{} Right to Work in UK \,\textbar\, Available to join Immediately
\end{center}

% ================= PROFILE =================
\needspace{5\baselineskip}\section{Profile}
Builds the modelling a brief needs and reports its limits plainly: gradient boosting (XGBoost churn), LSTM time series forecasting with time ordered validation, text clustering (SBERT, K-Means, DBSCAN) and 11,459 per gene models (train Pearson 0.807 vs test 0.545). Works with clients: scoped a loosely defined brief at 180 Degrees Consulting into a stakeholder demo (five star rating) and contributed to the client\textquotesingle{}s production MCP server, now in live use. Deploys on AWS and GCP (Lambda, Cloud Run) with CI/CD and MLflow experiment tracking. Owns data workflows end to end in Python and SQL, from merging structured, API and web scraped data through feature creation, tuning, validation and deployment (Docker, GitHub Actions CI/CD, MLflow, AWS Lambda).

% ================= SKILLS =================
\needspace{5\baselineskip}\section{Skills}
\skillgroup{Machine Learning, Deep Learning \& Statistics}{Predictive modelling (gradient boosting with XGBoost, random forest, scikit-learn), clustering (K-Means, HAC, DBSCAN, Leiden), deep learning (PyTorch, TensorFlow/Keras, LSTM, GRU, BiRNN, CNN), time series forecasting (cyclic time encoding, TimeSeriesSplit), anomaly detection (Isolation Forest), class imbalance handling (SMOTE), hyper parameter tuning (grid search, two stage Bayesian), model evaluation (held out, cross validation, Pearson, MCC), statistical inference and experimental design \textit{(confidence intervals, bootstrap, hypothesis testing; MSc coursework)}, MLflow experiment tracking}
\skillgroup{Python, Data Quality \& Cloud}{Python, pandas, NumPy, SciPy, Matplotlib, Seaborn, Plotly, SQL, PostgreSQL, R (tidymodels), identifier crosswalks, parquet, reproducible pipelines, AWS, GCP, Docker, Kubernetes (AWS EKS), serverless deployment (Lambda, Cloud Run), containerisation, GitHub Actions CI/CD, FastAPI, TensorFlow Serving, pytest, Git}
\skillgroup{NLP, LLMs, Embeddings \& Evaluation}{NLP (text preprocessing, TF-IDF, SBERT embeddings, text clustering), Hugging Face Transformers, large language models (Claude, OpenAI, open weight via Ollama), prompt engineering, RAG with embeddings (ChromaDB), agents (LangChain, LangGraph), team LoRA fine tune of Llama, LLM evaluation (hallucination rate analysis, deterministic guardrails)}
\skillgroup{Client Delivery \& Communication}{Scoping ambiguous problems with clients, stakeholder demos, working with client engineering teams, code review, transparent reporting of model limitations, root cause analysis, technical documentation}

% ================= EXPERIENCE =================
\needspace{5\baselineskip}\section{Work Experience}
\needspace{4\baselineskip}
\role{AstraZeneca}{Bristol, UK (Remote)}{May 2026 – Sep 2026}{Data Scientist}
\begin{itemize}
  \item Needed an LLM a research team could trust before it shipped: designed a deterministic hallucination guard architecture and fail safe data resolution via live Ensembl/HGNC APIs with a local \num{19,213 gene parquet fallback} and circuit breakers.
  \item Diagnosed where it struggled with a three tier, \textasciitilde{}65 test evaluation harness (tool level → API → live) and hallucination rate analysis, debugging failures layer by layer, and used the same harness to validate the Groq to Bedrock migration after a rate limit evaluation.
  \item Built the Python data pipeline: reconciled \num{14 heterogeneous source datasets} through identifier crosswalks so the agent used trusted data, then deployed it as a Docker image on AWS Lambda (ECR, least privilege IAM) and documented the architecture and evaluation.
\end{itemize}
\gap
\needspace{4\baselineskip}
\role{180 Degrees Consulting Bristol}{Bristol, UK (Remote)}{Sep 2025 – Dec 2025}{Data Consultant (AI Engineering \& AI System)}
\begin{itemize}
  \item A product startup client needed an agentic AI prototype turned into software their team could use: led the architecture for a five person cross functional team, designed hierarchical agent workflows (MCP, Exograph GraphQL, PostgreSQL for persistent storage and retrieval across agent sessions) and built the Python prototype (FastMCP, Ollama, GitHub Actions CI), then ran a live demo with client stakeholders \num{(five star rating from senior leadership)} within the three month engagement.
  \item Proposed the request flow: a user\textquotesingle{}s text goes through the API and a customised prompt to the MCP server, where one LLM instance writes GraphQL queries against the client\textquotesingle{}s Convex database and a second instance runs the statistical analysis before the result returns to the user.
  \item Carried it into production with the client\textquotesingle{}s engineering team: contributed to their production TypeScript/Node.js MCP server (OpenAI based query generation, Convex), now in live use, and reviewed pull requests in the client\textquotesingle{}s production repository alongside their CTO.
\end{itemize}
\gap
\needspace{4\baselineskip}
\role{Freelance}{Chennai, India}{Jan 2025 – Aug 2025}{AI/ML Engineer}
\begin{itemize}
  \item Built scikit-learn pipelines and LSTM/BiRNN/CNN models in TensorFlow/Keras, deployed as FastAPI services in Docker.
  \item Built NLP and agentic RAG workflows with Hugging Face Transformers, Ollama, LangChain and LangGraph.
  \item Designed reproducible scikit-learn pipelines with feature engineering (temporal features, interaction terms, target encoding), model selection, cross validation and hyperparameter tuning across classification and regression tasks.
\end{itemize}
\gap
\needspace{4\baselineskip}
\role{Avasoft}{Chennai, India}{Jan 2023 – Sep 2024}{Software Development Engineer (Jun 2023 – Sep 2024); Software Development Engineer Trainee (Jan 2023 – May 2023)}
\begin{itemize}
  \item Implemented offline first data persistence (Isar) with background synchronisation in the Flutter FinTech app, keeping financial records available in low connectivity conditions.
  \item Shipped features in a production fintech (financial services) application as part of a team: integrated Twelve Data (real time market feeds) and Plaid (bank linking, transaction imports) via REST APIs, with SSL pinning and AWS Cognito SSO.
  \item As a trainee, modelled user profile, transaction history and event booking data in Core Data schemas with CRUD and session state management, and automated a manual data collection process with a JavaScript web scraping pipeline that surfaced market trend signals used for pricing decisions.
\end{itemize}

% ================= PROJECTS =================
\needspace{5\baselineskip}\section{Projects}
\needspace{4\baselineskip}
\setbox0=\hbox{\textbf{Banking Customer Churn: propensity model}\hspace{1em}\url{https://github.com/Thiruvel-AP/Banking-Customer-Churn-Prediction}}%
\setbox2=\hbox{{\small\textit{R, tidymodels, XGBoost}}\hspace{1em}{\small\url{https://github.com/Thiruvel-AP/Banking-Customer-Churn-Prediction}}}%
\ifdim\wd0>\linewidth
\ifdim\wd2>\linewidth
\textbf{Banking Customer Churn: propensity model}\\ {\color{muted}\small\textit{R, tidymodels, XGBoost}}\\ \hspace*{\fill}{\small\url{https://github.com/Thiruvel-AP/Banking-Customer-Churn-Prediction}}\par
\else
\textbf{Banking Customer Churn: propensity model}\\ {\color{muted}\small\textit{R, tidymodels, XGBoost}}\hfill{\small\url{https://github.com/Thiruvel-AP/Banking-Customer-Churn-Prediction}}\par
\fi
\else
\project{Banking Customer Churn: propensity model}{R, tidymodels, XGBoost}{\url{https://github.com/Thiruvel-AP/Banking-Customer-Churn-Prediction}}
\fi
\begin{itemize}
  \item Predicted bank customer churn with a supervised XGBoost propensity classifier in R (tidymodels): two stage Bayesian hyperparameter tuning, SMOTE inside the preprocessing recipe and MCC evaluation.
\end{itemize}
\gap
\needspace{4\baselineskip}
\setbox0=\hbox{\textbf{Mining Insights from Customer Feedback: NLP \& unsupervised analytics (MSc team project)}\hspace{1em}\url{https://github.com/Thiruvel-AP/TextAnalytics}}%
\setbox2=\hbox{{\small\textit{Python, scikit-learn, SBERT, K-Means, HAC, DBSCAN}}\hspace{1em}{\small\url{https://github.com/Thiruvel-AP/TextAnalytics}}}%
\ifdim\wd0>\linewidth
\ifdim\wd2>\linewidth
\textbf{Mining Insights from Customer Feedback: NLP \& unsupervised analytics (MSc team project)}\\ {\color{muted}\small\textit{Python, scikit-learn, SBERT, K-Means, HAC, DBSCAN}}\\ \hspace*{\fill}{\small\url{https://github.com/Thiruvel-AP/TextAnalytics}}\par
\else
\textbf{Mining Insights from Customer Feedback: NLP \& unsupervised analytics (MSc team project)}\\ {\color{muted}\small\textit{Python, scikit-learn, SBERT, K-Means, HAC, DBSCAN}}\hfill{\small\url{https://github.com/Thiruvel-AP/TextAnalytics}}\par
\fi
\else
\project{Mining Insights from Customer Feedback: NLP \& unsupervised analytics (MSc team project)}{Python, scikit-learn, SBERT, K-Means, HAC, DBSCAN}{\url{https://github.com/Thiruvel-AP/TextAnalytics}}
\fi
\begin{itemize}
  \item Needed structure in unstructured airline support tickets: in a five person team, coordinated per member Git branches merged through pull requests and ran the team Trello board.
\end{itemize}

% ================= EDUCATION =================
\needspace{5\baselineskip}\section{Education}
\needspace{3\baselineskip}
\roleNoLoc{University of Bristol}{Sep 2025 – Sep 2026}{MSc Data Science, Distinction \textit{(Expected)}}
{\color{muted}\small Large scale Data Engineering, Statistical Computing \& Empirical Methods, AI \& Text Analytics, Visual Analytics, Data Science Methods \& Practices, Technology, Innovation, Business \& Society}\par
\gap
\needspace{3\baselineskip}
\roleNoLoc{St. Joseph\textquotesingle{}s College of Engineering}{Aug 2019 – Apr 2023}{B.E. Electrical and Electronics Engineering, 91\%, First Class with Distinction}

% ================= CERTIFICATIONS & ACHIEVEMENTS =================
\needspace{5\baselineskip}\section{Certifications \& Achievements}
\begin{itemize}
  \item Model Context Protocol: Advanced Topics, Anthropic, \textit{Jun 2026}
  \item Introduction to Model Context Protocol, Anthropic, \textit{Jun 2026}
  \item AWS Academy Cloud Foundations, \textit{Nov 2025}
  \item 180 Degrees Consulting Data Science Consultant Certificate, \textit{Dec 2025}
  \item Bristol PLUS Award, \textit{Mar 2026}
\end{itemize}

\end{document}